\selectlanguage{english}
\noindent\textbf{Relative homotopy types, local systems, and integration.}
These pages of Chapter VI pass from several models for homotopy types over a
fixed category to the operations that make those homotopy types into a
derivator. The intervening abstraction is substantial. Grothendieck first
organizes the comparison of fifteen categories over a base $S$, then isolates
an argument about fibrant replacement that applies in a general category with
weak equivalences, and finally returns to categories to construct inverse
image and homological direct image. The common concern is the independence
of a homotopy construction from its presentation. Different presentations
make different operations literal: a split fibration makes restriction
literal, whereas an arbitrary category over $S$ makes integration literal.
The comparison theorems let one use both conveniences in the same theory.

There are three localizations of $\Cat/S$, and keeping their roles distinct is
the main prerequisite for reading the chapter. For a functor $u:X\to Y$ over
$S$, write $W_S$ for weak equivalence of the total categories, and write
$W_S^{\mathrm d}$ and $W_S^{\mathrm g}$ for the conditions
\[
 s\backslash X\longrightarrow s\backslash Y\in W,
 \qquad
 X/s\longrightarrow Y/s\in W
 \quad(s\in S),
\]
respectively. Here $s\backslash X$ is the comma category of arrows
$s\to p_X(x)$; $X/s$ uses arrows $p_X(x)\to s$. These are comma categories,
not strict fibres. The first fine localization presents contravariant
homotopy coherent diagrams on $S$, the second presents the covariant
version, and the coarse localization presents locally constant objects.
Throughout the first theorem--scholium its strong fibration axiom L5 must be
read together with the axioms recalled at manuscript pages 43--45. The
edition points forward to a characterization of $W_\infty$. Its precise
form excludes the maximal localizer consisting of all functors: Chapter
XII, manuscript page 12, assumes detection of connected components, later
weakened to nondegeneracy. With that qualification the combination forces
$W=W_\infty$ \citep[Chapter XII, manuscript pp.~12 and 20--22]{Grothendieck1990XIIa}.
The general letter $W$ does not license extending every assertion to an
arbitrary homological localization. By contrast, the abstract second
theorem--scholium has its own explicitly stated hypotheses.

The locally constant category is a full reflective subcategory of the fine
one. Grothendieck constructs a replacement $X\to C(X)$ over $S$ whose target
is Hot-fibred and whose total-category map is a weak equivalence. It induces
\[
 \rho_S:\HOT_W(S)\rightleftarrows
 \HOT_W^{\mathrm{lc}}(S):i_S,
 \qquad \rho_S\dashv i_S,
 \qquad \rho_S i_S\simeq\id.
\]
The unit $X\to C(X)$ becomes invertible after coarse localization; it need
not do so after fine localization. The apparently small distinction is what
allows the fine theory to retain variation along noninvertible arrows of
$S$. For ordinary spaces one may interpret a locally constant object as a
family over the classifying space $BS$, with all its higher monodromy. Its
underlying functor on the fundamental groupoid is not a complete invariant.
Our example comparing the Hopf circle bundle with the trivial circle bundle
over a model of $S^2$ makes this loss of information explicit. The relevant
modern comparison is developed by \citet[\S\S2--4]{Cisinski2009LocallyConstant};
it is a comparison of coherent homotopy theories, not an identification with
ordinary functors into a homotopy category.

\medskip
\noindent\textbf{What replacement must accomplish.}
To prove the reflection, Grothendieck has to compare two ways of inserting a
trivial path:
\[
 i_{C(X)},\ C(i_X):C(X)\rightrightarrows C^2(X).
\]
They are not the same map of categories. Concatenation supplies a common
retraction, which becomes invertible in the relevant localization and makes
the maps equal there. This is an effective geometric proof of the second
triangle identity. The later abstraction separates the essential features:
a full subcategory of fibrant objects presents the coarse localization, and
coarse and fine weak equivalences agree between those fibrant objects.
Functorial replacement proves the first assertion directly. Without
functorial replacement, a lifting hypothesis and a factorization of the
relative diagonal prove that the induced arrow is independent of the chosen
lift. We give the naturality and independence arguments in full, including
the stability assumptions needed to keep the relative path object fibrant.

The commentary also extracts a simplification of the second
theorem--scholium. Under its functorial hypothesis A, the extra composition
of paths in A$'$ is unnecessary for the adjunction. Naturality equates the
two composites with $i_X$ in the coarse localization. Cancellation there,
followed by full faithfulness of the fibrant-subcategory presentation,
already equates $i_{C(X)}$ and $C(i_X)$ in the localization of the fibrant
subcategory. That equality maps to the fine localization because its weak
equivalences agree there. This is a formal consequence of the preceding
comparison theorem, proved in the development block; it is not a claim of
historical priority. A second development applies the same mechanism to
weak equivalences stable under every base change, answering the question at
manuscript page 84 under the stated fibration assumptions.

The historical setting explains the attention paid to this elementary
looking machinery. Section 69 of \emph{Pursuing Stacks} already connects
diagram categories with the inadequacy of a single derived category for
trace and determinant constructions, and proposes integration of relative
homotopy types by total categories
\citep[\S69]{Grothendieck1983}. Chapter VI works out the comparisons needed
to turn that proposal into operations on localized categories. The model
category route and the derivator route subsequently became complementary:
\citet[Theorem 6.11]{Cisinski2003} constructs homotopy right Kan extensions
from model categories with small limits, with the dual statement furnishing
left Kan extensions when small colimits exist. These results supply a
general existence theory beyond the particular geometric models used here.

\medskip
\noindent\textbf{Changing the base.}
For $f:S'\to S$, inverse image is first defined by precomposition of
$S^{\mathrm{op}}$-diagrams. On a smooth category over $S$ it is computed by
the strict pullback $X\times_S S'$. An arbitrary representative must first
be replaced by the split-fibration construction $\Phi_S(X)$:
\[
 f^*\hot_{W,S}(X)\simeq
 \hot_{W,S'}\bigl(\Phi_S(X)\times_S S'\bigr).
\]
The replacement $C_S(X)$ answers a different question: it first reflects
the object into the locally constant subcategory. Corollary 2 on manuscript
page 120 conflates these operations. Its assertion involving $C_S(X)$ is
false for a general $X$; the error is present in the edition. The walking
arrow gives a complete counterexample. We retain the printed assertion and
place the two correct formulas beside it. The same example also explains
why strict pullback need not preserve the fine weak equivalences on all of
$\Cat/S$.

Proper base changes provide another way to make strict pullback compute
the derived operation. Grothendieck presents the required preservation of
colocal equivalences as a conjecture at manuscript page 122, but its
annotation points to his own subsequent proof. The full edition confirms
that proof on manuscript pages 168--169. It is therefore not a currently
open conjecture, nor an omission that needs to be attributed to a later
author. Later systematic accounts express this calculus through asphericity
and homotopy exact base-change squares; see
\citet[Theorems 4.13 and 5.4]{Maltsiniotis2005Asphericite}. The converse
implication in Proposition 2 at manuscript pages 123--126 has a different
status. Grothendieck explicitly notices that invertibility after
localization need not imply membership in the specified weak-equivalence
class. Under strong saturation of the ambient $W$, the forgetful functor to
$\Cat[W^{-1}]$ supplies the missing implication. This includes the ordinary
weak homotopy equivalences.

\medskip
\noindent\textbf{Integration and the questions it leaves.}
Direct image is especially simple in the slice presentation. For a category
$X'\to S'$, compose its structural map with $f:S'\to S$. This operation
preserves fine equivalences because the local testing maps are stable under
base change. It consequently defines $f_!$ on relative homotopy types. The
adjunction $f_!\dashv f^*$ still requires care: after composing to $S$, a
previously fibrant representative need not remain fibrant. The unit uses
$X'\to\Phi_S(X')$, and for an arbitrary $X\to S$ the counit is the localized
composite
\[
 \Phi_S(X)\times_S S'\longrightarrow\Phi_S(X)
 \xrightarrow{i_X^{-1}}X.
\]
One must invert $i_X$ after projecting, rather than assume that its strict
pullback is invertible. Grothendieck himself rejects that false route. The
last displayed triangle of the supplied excerpt is completed by naturality
on the next manuscript page; our commentary gives both triangle identities
without mistaking the excerpt boundary for a gap.

Three questions here are particularly instructive today. First, localization
does commute with finite products of relative categories, but the analogous
statement for arbitrary infinite products is false without additional
hypotheses. Unbounded lengths of necessary zigzags produce an explicit
counterexample. Model presentations with a uniform description of derived
morphisms avoid that obstruction and give the required coproduct axiom for
the usual homotopy derivator. Second, the assignment
$S\mapsto\HOT^{\mathrm{lc}}(S)$ is not itself the ordinary homotopy derivator
on all categories. It has both adjoints to restriction, as proved by
\citet[Corollary 3.10]{Cisinski2009LocallyConstant}, but its pointwise Kan
extension axiom fails already for an endpoint inclusion into $[1]$. This
does not contradict the same paper's Theorem 3.12: that theorem fixes a
category $A$ and varies a separate diagram category $I$, forming the
derivator of locally constant $A$-objects. Keeping the base variable and the
diagram variable separate resolves the apparent contradiction.

Third, the cancelled assertion about Der~2 cannot be recovered from strong
saturation alone. We give a finite relative category with four objects and
fourteen arrows, whose weak equivalences are strongly saturated, but whose
localized arrow category contains a noninvertible endomorphism evaluated
as the identity at both vertices. Its noninvertibility is detected by an
explicit functor from the whole arrow category to sets, so no unproved
fully-faithfulness assertion about localization of a subcategory is needed.
The proof also isolates the positive hypothesis: a three-arrow
representation of localized diagram morphisms, together with strong
saturation, makes evaluation conservative. This agrees with the later
cofibration and model-category results
\citep[Theorem 3.33]{Lenz2018}. The counterexample is supplied as a
commentators' development; the literature check supports no claim of
priority.

\medskip
\noindent\textbf{Reading this companion.}
The audit records the status of the assertion at its location in the
manuscript and gives later completions separately. Edition slips remain in
the French text; verified transcription errors are corrected. References
to earlier commentary use the supplied DER\_I-a index, especially its
discussions of variance, mates, cofinality, and local constancy, without
repeating those discussions. The source PDF supplied with the package
covers edition pages 42--84, while the TeX covers pages 39--88. The omitted
margins were checked against the editors' complete PDF of the same
9 November 2001 edition: all 43 overlapping page content streams agree
byte for byte. This supplement also permits the limited forward checks
just identified; it does not enlarge the mathematical scope of the audit
beyond the supplied text.
